\begin{remark}\label{rl-000Z}\label{rmk:two-modifications-after-the-residue-independence-lemma-statement}
  Before proceeding, two remarks are in order. First, the same specialization map \(\sigma:A^T_*(\mathcal X)\to A^T_*(\mathfrak C)\) occurs in both \(\eta^!_1\) and \(\eta^!_2\), hence it suffices to prove that for every \(\beta\in A^T_*(\mathfrak C_{\mathcal Z/\mathcal X})\),
  \begin{align}\label{eqn:independence-lemma-simplification}
    e_T(\mathfrak E_1)^{-1}\cap 0^!_{\mathfrak E_1}\big((j_1)_*\beta\big) = e_T(\mathfrak E_2)^{-1}\cap 0^!_{\mathfrak E_2}\big((j_2)_*\beta\big).
  \end{align}
  This removes the stack \(\mathcal X\) and the map \(\sigma\) from the discussion entirely.

  Second, once the specialization map \(\sigma\) is removed, the cone structure of \(\mathfrak C\) is no longer relevant, and we can simply take \(\mathfrak C\) to be any algebraic stack \(T\)-equivariant over \(\mathcal Z\) admitting a \(T\)-equivariant closed embedding into \(\mathfrak E\). This is not a superfluous point; the proof of Lemma~\ref{lem:independence-of-Res-map-from-choice} requires this generality.
\end{remark}
